% Cross-modality study (v4): the fourth analysis plan SAP_v4.md and its design DESIGN_v4.md, hash-frozen in freeze_v4.json on 8 November 2026
% before the first v4 training run; later changes are in DEVIATIONS.md (Appendix app:deviations). Tables: paper/tables/v4_*.tex, written by
% tools/xm_tables.py from results/final_xm_stats.json (tools/xm_stats.py), results/v4/delong.json, results/v4/pointing_summary.json and
% results/v4/uq_cls_*.json; figures: paper/figures/fig_v4_*.pdf (tools/fig_v4.py). Numbers in the text are copied from those files.
\section{Question}
\label{sec:xm_intro}
The earlier chapters used CT in one way only: to render training images for a bone suppressor. No classifier ever saw a CT volume. This
chapter asks whether CT can teach a chest radiograph classifier directly. CT resolves in three dimensions what a radiograph superimposes,
so a classifier trained on CT learns its findings without the overlap of ribs, heart and mediastinum. If that knowledge passes to an X-ray
model through images of the same patients, the X-ray model might classify real radiographs better, look at the right place and report
fewer findings that are not there.

The chapter builds and tests such a framework: an explainable and hallucination-aware Swin transformer for cross-modality chest CT and
X-ray classification. One Swin family, SwinV2-T \citep{liu2022swinv2}, serves both modalities. A CT teacher learns six findings from
volumes of CT-RATE \citep{hamamci2026ctrate}. An X-ray model learns the same findings from frontal projections of the same volumes, once on
its own (P0) and once with knowledge distillation from the teacher (P1) \citep{hinton2015distilling}. Both then initialise classifiers of
the 14 NIH ChestX-ray14 diseases \citep{wang2017chestxray}, beside a classifier started from ImageNet weights. Five hypotheses were fixed
before any v4 model was trained, in a fourth analysis plan that is pre-specified and hash-frozen (8 November 2026) but not deposited in a
public registry; they continue the numbering of the earlier plans. H11 (primary) and H12 ask whether distillation from CT improves the
classification of real radiographs; H13 compares the Swin classifier with the DenseNet-121 of Chapter~\ref{ch:classifier}; H14 asks
whether the CT and the X-ray models attend to the same place in the same patient; and H15 asks whether the CT-distilled classifier
hallucinates less when an image gains anatomy but no disease.

\section{Materials and methods}
\label{sec:xm_methods}

\subsection{Data}
\label{sec:xm_data}
\paragraph{Paired CT and projections.} The Qure.ai ct2xr projections \citep{angaitkar2026anatomy} are frontal projections of CT-RATE
volumes, one per volume, each with six findings taken from its scan: lung opacity, pleural effusion, atelectasis, cardiomegaly,
consolidation and lung lesion. They were split by patient as in Chapter~\ref{ch:v2} (\num{21051} training, 406 validation and 430 test
projections). The CT subset holds one reconstruction per patient: every validation and test patient of that split and \num{1600} training
patients drawn with a fixed seed, \num{2156} volumes in all (\num{1600} training, 259 validation, 297 test). Volumes that CT-RATE lists as
not showing the chest and volumes not scanned head first and supine were never selected, and their projections were left out of the
projection models as well, which leaves \num{20533} training, 396 validation and 415 test projections. Each volume was reduced inside the
cloud kernel to 24 axial slices spread over the lung extent, with three Hounsfield windows as channels (lung, mediastinum, bone) and
cropped to the body at 256\,px; the lung box and an anteroposterior attenuation projection were kept for the heat maps. CT-RATE is
released under CC BY-NC-SA~4.0; following the design, nothing derived from it left the cloud accounts except metrics and the
$32\times32$ heat-map frames, with the logged exceptions of Appendix~\ref{app:deviations} (teacher outputs, predictions and the weights of
the distilled models, which stayed private on the author's computer).

\paragraph{Real radiographs.} NIH ChestX-ray14 and VinBigData were used with the splits of Chapter~\ref{ch:classifier}
(Section~\ref{sec:cls_data}): NIH \num{77890} training, \num{8634} validation and \num{25596} test images with the 984 expert boxes of the
test list; VinBigData \num{2943} test images, used for zero-shot transfer only. A fixed label map links the six projection findings to
the real data sets: pleural effusion, atelectasis, cardiomegaly and consolidation map to the classes of the same name, and lung lesion to
NIH Mass or Nodule (VinBigData Nodule/Mass). These five are the exactly mapped shared findings of H11 and H12. Lung opacity maps exactly to
VinBigData Lung Opacity but only approximately to NIH Infiltration; it is reported but not tested.

\subsection{Models and training}
\label{sec:xm_models}
All models use SwinV2-T (\texttt{swinv2\_tiny\_window8\_256}, ImageNet-1k weights from timm \citep{wightman2019timm}) and the class-query
decoder of Chapter~\ref{ch:classifier} (Section~\ref{sec:cls_methods}); learning rate $3\times10^{-4}$ with the encoder at 0.25 of it, and
five seeds per model (Fig.~\ref{fig:xm_design}).

\begin{figure}[htbp]
\centering
\safefig{figures/fig_v4_design.pdf}
\caption{Study design. A 2.5D CT teacher learns six findings on the CT subset; projection models learn them on the Qure.ai projections
without (P0) and with (P1) distillation from the teacher; NIH classifiers start from ImageNet (F$_\mathrm{in}$), P0 (F$_\mathrm{p0}$) or P1
(F$_\mathrm{p1}$). All models are SwinV2-T with the class-query decoder, five seeds each; the DenseNet-121 of
Chapter~\ref{ch:classifier} is the CNN reference.}
\label{fig:xm_design}
\end{figure}

\paragraph{CT teacher.} A 2.5D network: the X-ray encoder runs on each of the 24 slices with the three windows as colour channels; every
slice yields $8\times8$ tokens with a two-dimensional sine code and a learned slice-depth embedding, and the class-query decoder attends
over the tokens of all slices. A full 3D Swin did not fit a T4 GPU at a useful batch size, and the 2.5D design keeps the ImageNet weights
and the architecture of the X-ray model. It was trained for 20 epochs with four volumes per step, keeping the best validation epoch.

\paragraph{Projection models.} P0 learns the six findings on the projections for 15 epochs (batch 32). P1 is trained the same way but,
on the \num{1600} training rows whose volume is in the CT subset, also matches the teacher's logits (binary distillation at temperature
2, weight 1) and aligns its pooled decoder output with the teacher's through a linear head (cosine loss, weight 0.5); half of every batch
is drawn from these paired rows. P1 seed $k$ distils from the teacher of slot $k$ (below).

\paragraph{NIH classifiers.} F$_\mathrm{in}$ starts from ImageNet, F$_\mathrm{p0}$ from P0 and F$_\mathrm{p1}$ from P1 (encoder and
decoder loaded, class heads new); each is trained on the 14 NIH diseases for 10 epochs (batch 32) with the radiograph as its only input,
keeping the best validation epoch. The five DenseNet-121 radiograph-only runs of Chapter~\ref{ch:classifier} are the CNN reference; they
were trained for 20 epochs at the full encoder learning rate, the pre-specified rule for batch-normalised encoders, so H13 compares
encoders trained with different budgets, neither tuned further.

\paragraph{Seed handling.} CT teacher seed 2 did not learn: its validation macro-AUROC stayed between 0.46 and 0.54 and its logits were
constant over all volumes. Before any P1 model of seed 2 was trained, a rule was fixed and logged: train seeds 5 and 6 with the unchanged
recipe and give slot 2 to seed 5 unless its logits were constant (per-class standard deviation below 0.01 over the training and validation
volumes in every class), otherwise to seed 6. Seed 5 learned and took slot 2; seed 6 was not used. The teachers of slots 0--4 are
therefore seeds 0, 1, 5, 3 and 4.

\subsection{Explainability}
\label{sec:xm_xai}
\paragraph{X-ray.} Grad-CAM++ on the last encoder stage \citep{chattopadhay2018gradcampp}, Integrated Gradients
\citep{sundararajan2017axiomatic} and the attention of the last decoder layer were computed as in Section~\ref{sec:cls_xai}, and the
pointing game was played on the 984 NIH boxes with the same tolerance of 15\,px at 256\,px \citep{zhang2018topdown}, for each seed of
F$_\mathrm{in}$, F$_\mathrm{p0}$, F$_\mathrm{p1}$ and the DenseNet-121 reference.

\paragraph{CT and cross-modal agreement (H14).} For each CT subset test scan positive for a finding, Grad-CAM++ was computed on every slice
of the teacher, the slice maps were stacked into a saliency volume and summed along the anteroposterior axis, which gives a frontal map in
the geometry of the paired projection. The CT map and the X-ray map of the paired projection (P1, and P0 for comparison) were cropped to
their own lung boxes, resized to $32\times32$, scaled to sum 1 and averaged over the five seed pairs (teacher of slot $k$ with X-ray seed
$k$). Agreement is the Pearson correlation of the matched maps; its null pairs the CT map with the X-ray maps of the other scans positive
for the same finding. The pre-specified orientation check used the frames alone, never a heat map: the CT attenuation frames correlated
better with the projections as they were (0.472) than mirrored (0.435), so no X-ray map was mirrored.

\subsection{Hallucination tests}
\label{sec:xm_halluc}
A classifier hallucinates when it reports a finding the image does not contain, or misses one it does. The tests use the
\num{9861} NIH test images without any finding and each model's seed-mean probabilities, with a per-class threshold at 90\,\% specificity
on the validation split. \emph{Finding-free perturbations:} rib emphasis (the radiograph plus half of the bone image of the suppressor of
Chapter~\ref{ch:v2}, which adds anatomy but no disease), Gaussian noise ($\sigma=0.02$) and a gamma of 0.8; the hallucination rate is the
share of images in which some class moves from at most to above its threshold. \emph{Counterfactual lesions:} a nodule from the lesion
library of Chapter~\ref{ch:audit} is pasted into a lung field of images whose Mass and Nodule scores are at most their thresholds; sensitivity
is the share in which either score crosses its threshold, and the score change is reported by nodule-size tertile. \emph{Off-target
attribution:} the share of Grad-CAM++ mass outside the lung field for positive predictions of lung findings. Seed groups of one model run on
different accounts were combined first (probabilities averaged with the number of seeds as weights, thresholds from the averaged
validation probabilities).

\subsection{Statistics}
\label{sec:xm_stats}
A model's prediction is the mean of its seeds' probabilities. H11 compares F$_\mathrm{p1}$ with F$_\mathrm{in}$ in the mean test AUROC of
the five exactly mapped shared findings; H12 makes the same comparison for F$_\mathrm{p1}$ against F$_\mathrm{p0}$ and for F$_\mathrm{p0}$
against F$_\mathrm{in}$; H13 compares F$_\mathrm{in}$ with the DenseNet-121 reference in the macro-AUROC over all 14 NIH diseases. All
three use a paired patient-cluster bootstrap (\num{10000} resamples) with two-sided percentile $p$-values (the $+1$ rule); H11 and the two
H12 contrasts form one Holm family. H14 is the mean per-map difference between matched and mismatched correlation with a scan-cluster
bootstrap, and H15 the mean per-image difference in the rib-emphasis flip between F$_\mathrm{p1}$ and F$_\mathrm{in}$ with a
patient-cluster bootstrap. The plan lists as descriptive: the CT teacher, the projection models and their zero-shot transfer to NIH and
VinBigData, per-class DeLong tests \citep{delong1988auc}, conformal risk control of the false-negative rate at 0.10
\citep{angelopoulos2022crc} for the Swin classifiers (calibrated on the NIH validation predictions as in Section~\ref{sec:cls_uq}) and the
pointing game. How the DeLong tests are grouped (Holm over the 15 shared-finding tests and over the 14 classes of the H13 contrast) and how
the pointing game is summarised (hit rate of the seed-mean hits with an image-cluster bootstrap interval) were fixed after the hypothesis
results were known and are logged as such.

\section{Results}
\label{sec:xm_results}
All 32 training runs finished (seven CT teachers, ten projection models and 15 NIH classifiers), as did the 20 zero-shot transfers, the
ten cross-modal map runs, seven hallucination test runs and 20 pointing-game runs. One zero-shot run first failed because it started
before its model had been written and was rerun unchanged. Table~\ref{tab:xm_ct_seeds} and Fig.~\ref{fig:xm_ct_training} show the seven
CT teachers.

\begin{table}[htbp]
\centering\small
\caption{Every CT teacher run (unchanged recipe, 20 epochs): best validation and test macro-AUROC over the six findings of the CT subset.
Slot $k$ is the teacher of P1 seed $k$; seed 2 collapsed and seed 5 replaced it by the rule fixed beforehand.}
\label{tab:xm_ct_seeds}
\fittable{tables/v4_ct_seeds.tex}
\end{table}

\begin{figure}[htbp]
\centering
\safefig[88mm]{figures/fig_v4_ct_training.pdf}
\caption{Validation macro-AUROC of the CT teacher per epoch for all seven seeds. Seed 2 never left chance level, seed 6 learned late and
was not used, and seed 0 stalled near 0.60 for 12 epochs before it learned.}
\label{fig:xm_ct_training}
\end{figure}

\subsection{CT teacher and projection models}
\label{sec:xm_res_transfer}
The CT teacher reached a seed-ensemble test macro-AUROC of 0.821 on the 297 CT subset test volumes (Table~\ref{tab:xm_transfer}). It was
strong for cardiomegaly (0.951) and pleural effusion (0.942) and weak for lung lesions (0.608). The recipe was unstable across seeds: the
single teachers ranged from 0.737 to 0.822 among the five used, and two of seven runs collapsed or nearly did.

On the projections, distillation lowered accuracy slightly. The seed ensemble of P1 reached 0.806 against 0.814 for P0, and P1 was below
P0 in every seed (Table~\ref{tab:xm_seeds}, Fig.~\ref{fig:xm_transfer}a). Without any training on real radiographs, the projection
models transferred moderately to VinBigData (P0 0.739, P1 0.732) and weakly to NIH (P0 0.648, P1 0.640); distillation again gave the
lower value in both.

\begin{table}[htbp]
\centering\small
\caption{CT teacher, projection models and their zero-shot transfer: test AUROC of the seed ensembles over the six findings, with the
macro-average (descriptive). NIH lung opacity is the approximate mapping to Infiltration.}
\label{tab:xm_transfer}
\fittable{tables/v4_transfer.tex}
\end{table}

\begin{table}[htbp]
\centering\small
\caption{Test macro-AUROC of every single model (one seed each); the other tables use the seed ensembles. Projection models and zero-shot
transfer: six findings; NIH classifiers: 14 diseases.}
\label{tab:xm_seeds}
\fittable{tables/v4_seeds.tex}
\end{table}

\begin{figure}[htbp]
\centering
\safefig{figures/fig_v4_transfer.pdf}
\caption{Macro-AUROC of single seeds (open circles) and seed ensembles (diamonds). (a) Projection test split; (b, c) zero-shot transfer
of the projection models to NIH and VinBigData over the six shared findings; (d) NIH classifiers over the 14 diseases.}
\label{fig:xm_transfer}
\end{figure}

\subsection{Distillation and real radiographs (H11, H12)}
\label{sec:xm_res_h11}
Distillation from CT did not improve the classification of real radiographs (Table~\ref{tab:xm_hyp}, Fig.~\ref{fig:xm_forest}a). On the
five shared findings F$_\mathrm{p1}$ differed from F$_\mathrm{in}$ by $-0.0001$ $[-0.0016, 0.0014]$ (Holm $p=0.912$), so a gain larger
than 0.0014 is not compatible with the data. F$_\mathrm{p1}$ did not differ from F$_\mathrm{p0}$ ($-0.0008$ $[-0.0022, 0.0006]$, Holm
$p=0.742$), nor F$_\mathrm{p0}$ from F$_\mathrm{in}$ ($+0.0007$ $[-0.0008, 0.0023]$, Holm $p=0.742$). Over all 14 diseases the three Swin
classifiers reached 0.819 (F$_\mathrm{in}$), 0.822 (F$_\mathrm{p0}$) and 0.821 (F$_\mathrm{p1}$) (Table~\ref{tab:xm_macro}). Per shared
finding no difference survived the correction; the largest were $+0.0044$ for cardiomegaly with F$_\mathrm{p0}$ against F$_\mathrm{in}$
(Holm $p=0.093$) and $-0.0042$ for consolidation with F$_\mathrm{p1}$ against F$_\mathrm{in}$ (Holm $p=0.122$; Table~\ref{tab:xm_shared}).

\begin{table}[htbp]
\centering\footnotesize
\caption{Pre-specified hypotheses: differences of the seed ensembles with 95\,\% bootstrap intervals (\num{10000} resamples). Holm: H11
with the two H12 contrasts. The P0 row of H14 is descriptive.}
\label{tab:xm_hyp}
\fittable{tables/v4_hypotheses.tex}
\end{table}

\begin{figure}[htbp]
\centering
\safefig{figures/fig_v4_forest.pdf}
\caption{Pre-specified contrasts with 95\,\% bootstrap intervals. (a) AUROC contrasts of H11--H13; (b) H14, matched minus mismatched
correlation of the CT and X-ray heat maps; (c) H15, difference in the share of finding-free images flipped by rib emphasis. Filled
markers: $p<0.05$ (Holm-adjusted for H11 and H12).}
\label{fig:xm_forest}
\end{figure}

\begin{table}[htbp]
\centering\small
\caption{NIH test macro-AUROC of the seed ensembles over the 14 diseases and mean AUROC of the five shared findings (descriptive).}
\label{tab:xm_macro}
\fittable{tables/v4_macro.tex}
\end{table}

\begin{table}[htbp]
\centering\footnotesize
\caption{Shared findings on the NIH test split: AUROC of the seed ensembles and paired DeLong differences, Holm over the 15 tests
(descriptive).}
\label{tab:xm_shared}
\fittable{tables/v4_shared_delong.tex}
\end{table}

\subsection{Swin against DenseNet-121 (H13)}
\label{sec:xm_res_h13}
The DenseNet-121 reference was better than F$_\mathrm{in}$ over the 14 diseases by $0.0049$ $[0.0015, 0.0085]$ ($p=0.005$; the difference
F$_\mathrm{in}$ minus DenseNet-121 is $-0.0049$). The per-class tests, which were not planned in this form, locate almost all of it in one
class (Table~\ref{tab:xm_per_class}, Fig.~\ref{fig:xm_per_class}): for hernia, the rarest NIH class (86 positive test images),
DenseNet-121 reached 0.901 against
0.833 ($-0.0679$, Holm $p=0.0008$), which alone accounts for $-0.0049$ of the 14-class mean. The only other class that survived the
correction was infiltration ($-0.0062$, Holm $p=0.003$); the remaining 12 classes differed by less than 0.007 in either direction. On the
classes other than hernia the two encoders were practically equivalent.

\begin{table}[htbp]
\centering\footnotesize
\caption{NIH per-class test AUROC of the seed ensembles, with the paired DeLong difference F$_\mathrm{in}$ minus DenseNet-121 and Holm over
the 14 classes (descriptive).}
\label{tab:xm_per_class}
\fittable{tables/v4_per_class.tex}
\end{table}

\begin{figure}[htbp]
\centering
\safefig{figures/fig_v4_per_class.pdf}
\caption{(a) NIH per-class test AUROC of the seed ensembles. (b) F$_\mathrm{in}$ minus DenseNet-121 per class with 95\,\% DeLong
intervals; filled markers: Holm $p<0.05$ over the 14 classes. Hernia carries almost all of the H13 difference.}
\label{fig:xm_per_class}
\end{figure}

\subsection{Cross-modal agreement (H14)}
\label{sec:xm_res_h14}
The CT and X-ray heat maps did not agree. Over 419 maps of 227 CT subset test scans, the matched CT and P1 maps correlated at $-0.052$ on
average and the mismatched ones at $-0.060$, a difference of $+0.0076$ $[-0.0068, 0.0223]$ ($p=0.301$; Table~\ref{tab:xm_hyp},
Fig.~\ref{fig:xm_forest}b). With P0 the matched maps correlated less than the mismatched ones ($-0.0087$ $[-0.0228, 0.0056]$). Per
finding all mean correlations were negative, from $-0.008$ to $-0.140$, and the matched-minus-mismatched differences changed sign from one
finding to the next (Table~\ref{tab:xm_crossmodal}, Fig.~\ref{fig:xm_crossmodal}). The mean maps show why (Fig.~\ref{fig:xm_maps}). The
projected CT maps vary almost only from top to bottom, in horizontal bands that follow the importance of individual slices, while the
X-ray maps of most findings concentrate along the edges of the lung box; for cardiomegaly they peak near the centre, where the heart is,
and for lung lesions they are broad and weighted to the upper zones. Maps with such different structure are unlikely to correlate
positively, whatever the patient.

\begin{table}[htbp]
\centering\small
\caption{Cross-modal agreement per finding: mean Pearson correlation of the projected CT map with the X-ray map of the same scan (matched)
and with the X-ray maps of the other scans positive for the same finding (mismatched), maps averaged over the five seed pairs.}
\label{tab:xm_crossmodal}
\fittable{tables/v4_crossmodal.tex}
\end{table}

\begin{figure}[htbp]
\centering
\safefig[88mm]{figures/fig_v4_crossmodal.pdf}
\caption{Mean correlation of the CT map with the X-ray map of the same scan (filled) and of other scans positive for the same finding
(open), for P1 and P0; in parentheses the number of maps.}
\label{fig:xm_crossmodal}
\end{figure}

\begin{figure}[htbp]
\centering
\safefig{figures/fig_v4_maps.pdf}
\caption{Mean heat maps over the CT subset test scans positive for each finding: the CT teacher's Grad-CAM++ summed along the
anteroposterior axis (top) and the Grad-CAM++ maps of P1 and P0 on the paired projections, all cropped to the lung box at $32\times32$ and
scaled to sum 1 before averaging; each panel has its own colour scale. Only averages over scans are shown.}
\label{fig:xm_maps}
\end{figure}

\subsection{Hallucination tests (H15)}
\label{sec:xm_res_h15}
Rib emphasis flipped 13.7\,\% of the finding-free NIH test images for F$_\mathrm{in}$ and 12.9\,\% for F$_\mathrm{p1}$, a difference of
$-0.0080$ $[-0.0159, -0.0003]$ ($p=0.044$; Table~\ref{tab:xm_hyp}, Fig.~\ref{fig:xm_forest}c). H15 is secondary and was not corrected for
multiplicity. The comparison it was built for does not single out distillation: F$_\mathrm{p0}$, trained on the same projections without
the teacher, flipped 12.7\,\% and DenseNet-121 12.8\,\% (Table~\ref{tab:xm_halluc}, Fig.~\ref{fig:xm_halluc}). The clearest difference
between the models came from noise: it flipped 78.1\,\% of the images for DenseNet-121 and 27.3--29.3\,\% for the Swin classifiers, while a
gamma change flipped about 6\,\% for all four. Inserted nodules were rarely detected (sensitivity 2.8--3.2\,\%); the score change grew
with nodule size, from practically zero for nodules of 4.0--6.2\,mm to between $+0.015$ and $+0.019$ for those of 8.5--37.6\,mm
(Table~\ref{tab:xm_nodule}). Between 62\,\% (DenseNet-121) and 73\,\% (F$_\mathrm{p0}$) of the Grad-CAM++ mass of positive lung-finding
predictions fell outside the lung fields.

\begin{table}[htbp]
\centering\small
\caption{Hallucination tests on the \num{9861} finding-free NIH test images: share with some class above its threshold on the clean image,
share flipped by each perturbation, inserted-nodule sensitivity and the share of Grad-CAM++ mass outside the lung fields.}
\label{tab:xm_halluc}
\fittable{tables/v4_halluc.tex}
\end{table}

\begin{table}[htbp]
\centering\small
\caption{Counterfactual nodules: mean change of the Mass/Nodule score after insertion, by nodule-size tertile.}
\label{tab:xm_nodule}
\fittable{tables/v4_nodule_size.tex}
\end{table}

\begin{figure}[htbp]
\centering
\safefig{figures/fig_v4_halluc.pdf}
\caption{Hallucination tests. (a) Share of finding-free NIH test images in which some class crosses its threshold after a perturbation
that adds no disease. (b) Mean change of the Mass/Nodule score after inserting a nodule, by size tertile. (c) Share of Grad-CAM++ mass
outside the lung fields for positive lung-finding predictions.}
\label{fig:xm_halluc}
\end{figure}

\subsection{Pointing game}
\label{sec:xm_res_pointing}
The Swin classifiers' Grad-CAM++ maps rarely pointed into the expert boxes: 0.185 of the 984 targets for F$_\mathrm{in}$, 0.151 for
F$_\mathrm{p0}$ and 0.154 for F$_\mathrm{p1}$, against 0.385 for DenseNet-121 (Table~\ref{tab:xm_pointing}, Fig.~\ref{fig:xm_pointing}).
The decoder attention (0.479--0.572) and Integrated Gradients (0.615--0.645) did much better and differed little between the models. All
maps hit focal boxes (masses and nodules) less often than the other boxes. Grad-CAM++ of the Swin models often spreads along the image
borders (Fig.~\ref{fig:xm_xai}), as in the cross-modal maps.

\begin{table}[htbp]
\centering\small
\caption{Pointing game on the 984 NIH expert boxes (164 focal: Mass, Nodule): hit rate averaged over the five seeds, with a 95\,\%
image-cluster bootstrap interval (descriptive).}
\label{tab:xm_pointing}
\fittable{tables/v4_pointing.tex}
\end{table}

\begin{figure}[htbp]
\centering
\safefig{figures/fig_v4_pointing.pdf}
\caption{Pointing game. (a) Hit rate over all 984 boxes with 95\,\% image-cluster bootstrap intervals. (b) Focal (filled) and other
(open) boxes.}
\label{fig:xm_pointing}
\end{figure}

\begin{figure}[htbp]
\centering
\safefig{figures/fig_v4_xai_examples.pdf}
\caption{Heat maps of the seed-0 models on NIH test images with expert boxes (white): the first image of each box class among the 24
targets whose maps were saved; these 24 contain no mass or nodule box.}
\label{fig:xm_xai}
\end{figure}

\subsection{Uncertainty}
\label{sec:xm_res_uq}
Conformal risk control held the false-negative rate below its target for all three Swin classifiers: calibrated on the \num{8634} NIH
validation images, the test rate was 0.068--0.070 against a nominal 0.10, at the price of large label sets (6.33--6.45 of the 14 classes
per image; Table~\ref{tab:xm_conformal}, Fig.~\ref{fig:xm_conformal}). Selective classification by ensemble entropy behaved identically
for the three (area under the risk-coverage curve 0.0358).

\begin{table}[htbp]
\centering\small
\caption{Conformal risk control of the per-image false-negative rate at 0.10 on the NIH test split and selective classification by
$1-$entropy (descriptive).}
\label{tab:xm_conformal}
\fittable{tables/v4_conformal.tex}
\end{table}

\begin{figure}[htbp]
\centering
\safefig{figures/fig_v4_conformal.pdf}
\caption{(a) Risk-coverage curves of the Swin classifiers (Brier score of the images kept, most confident first). (b) Test false-negative
rate under conformal risk control at 0.10 (dashed); numbers: mean label-set size.}
\label{fig:xm_conformal}
\end{figure}

\section{Discussion}
\label{sec:xm_discussion}
Distillation from a CT teacher did not make a better radiograph classifier. On the shared findings the CT-distilled classifier matched
the one started from ImageNet within $\pm0.0016$ AUROC, and it matched the classifier pretrained on the same projections without the
teacher; on the projections themselves and in zero-shot transfer distillation was slightly worse. With five seeds, seed ensembles and a
plan fixed before training, the null result is informative: the interval excludes the gains that would justify a CT branch for
radiograph classification with this design.

Three observations suggest why. First, the teacher had little to give. On the six findings it was hardly better on CT (0.821) than P0 was
on the projections of the same test patients (0.814), and it was weakest exactly where CT should help most, lung lesions (0.608). Second, the
recipe was fragile: one of seven teachers collapsed, another nearly did, and the five used ranged from 0.737 to 0.822, so P1 seeds learned
from teachers of very different quality. Third, whatever P1 gained on the projections had to survive the move to real radiographs; the
zero-shot results (0.640--0.648 on NIH) show how wide that gap is, and the ten epochs of NIH fine-tuning may simply overwrite the
initialisation, as the equal results of F$_\mathrm{p0}$ and F$_\mathrm{in}$ suggest.

The explanation half of the framework failed in an instructive way. H14 asked whether the two modalities attend to the same place; they do
not even produce maps of the same kind. The 2.5D teacher's Grad-CAM++ mostly ranks slices, so its projected map varies from top to bottom
but hardly across, and the last stage of SwinV2-T gives an $8\times8$ grid whose Grad-CAM++ spreads along the image borders. The pointing
game shows the same weakness on real radiographs: Grad-CAM++ hit the boxes less than half as often for Swin as for DenseNet-121, while
decoder attention and Integrated Gradients localised about as well as for the CNN. For Swin classifiers the latter two are the maps to
report.

The hallucination tests separate robustness from accuracy. The Swin classifiers were not more accurate than DenseNet-121, but they were
far less sensitive to image noise (flip rate about 0.28 against 0.78), a property the accuracy numbers do not show. Under rib emphasis the
CT-distilled classifier flipped fewer images than the ImageNet-started one ($p=0.044$, uncorrected), but the classifier pretrained on the
projections without the teacher did the same, so the effect, if real, belongs to pretraining on projections rather than to CT. No model
reacted to inserted nodules of the two smaller size tertiles (up to 8.5\,mm), and even the largest ones moved the scores by less than
0.02.

\paragraph{Limitations.} The CT recipe was seed-unstable, and the replacement of seed 2 by seed 5 followed a rule fixed beforehand but
used a teacher trained later. The CT subset is small (\num{1600} training volumes) and its labels come from reports, as do those of the
projections. The NIH test split had been used in the earlier chapters with other models; H13's localisation in hernia comes from the
per-class tests, which were planned only as descriptive, and hernia has few positive test images. The cross-modal agreement depends on
Grad-CAM++ for both modalities, the weakest of the three maps on real radiographs; agreement of Integrated Gradients or attention maps was
not part of the plan. The example heat maps come from the first 24 saved targets, which contain no mass or nodule. VinBigData fine-tuning
(about 50 GPU hours) and the design's idea of using CT--X-ray disagreement as an abstention signal were not carried out. The teacher
outputs, the predictions and the weights of the distilled models left the cloud accounts for the author's computer, against the data
rule of the design; this is logged, and nothing derived from CT-RATE was released.

\section{Summary}
\label{sec:xm_summary}
A Swin framework for chest CT and X-ray was built and tested under a pre-specified, hash-frozen plan. Distilling a CT teacher into an X-ray
model did not improve the classification of real radiographs (H11, H12), the Swin classifier matched DenseNet-121 except in one rare class
(H13), and the CT and X-ray heat maps did not agree (H14). The Swin classifiers were much more robust to noise than DenseNet-121. Under
rib emphasis the CT-distilled classifier flipped slightly fewer finding-free images than the one started from ImageNet (H15, uncorrected),
as did the classifier pretrained on the projections without the teacher. Grad-CAM++ is a poor explanation for Swin classifiers;
Integrated Gradients and decoder attention localise about as well as for a CNN.
